%! Absolute Value Equations \& Inequalities — Unit 2, Lesson 2.2 — Warm-Up
\documentclass[10pt]{article}
\usepackage{algebra2-article}
\usepackage{algebra2-boxes}
\newcommand{\TallMath}[1]{$\displaystyle #1\rule[-1.4em]{0pt}{3.2em}$}
\newcommand{\numline}[2]{%
  \draw[->, semithick, charcoal] (#1-0.4,0)--(#2+0.4,0);
  \foreach \x in {#1,...,#2}{\draw[charcoal] (\x,0.10)--(\x,-0.10) node[below, font=\scriptsize]{$\x$};}}

\begin{document}

\pageheader{Unit 2, Lesson 2.2}{Warm-Up}
\namedateperiod

\vspace{0.10in}

\begin{notesbox}{Getting Ready: Distance, Solving, and Number Lines}
\small Today an absolute value asks a \emph{distance} question, and you'll solve the linear equation
inside the bars \emph{twice}. These three skills are the tools you'll need.

\vspace{4pt}
\textbf{1. Absolute value is distance from zero.} Evaluate each.
\[
  |-7| = \blank{1.2cm} \qquad |4| = \blank{1.2cm} \qquad
  |{-3}| + |5| = \blank{1.2cm} \qquad -|{-6}| = \blank{1.2cm}
\]
How many different numbers have an absolute value of $7$? \blank{2.2cm}\; Name them: \blank{2.4cm}

\vspace{6pt}
\textbf{2. Solve the two-step equation.} (You'll repeat this move once the bars split into two cases.)
\[
  2x - 3 = 7 \qquad\Rightarrow\qquad x = \blank{2.0cm}
  \hspace{1.2cm}
  2x - 3 = -7 \qquad\Rightarrow\qquad x = \blank{2.0cm}
\]

\vspace{6pt}
\textbf{3. Graph on a number line, then name it in interval notation.} Shade $x \ge -1$.
\begin{center}
\begin{tikzpicture}[scale=0.55]
  \numline{-5}{5}
\end{tikzpicture}
\end{center}
Interval notation: \blank{3.0cm} \qquad {\footnotesize (Use $[\,$ or $]$ for ``or equal to,'' $($ or $)$ otherwise.)}
\end{notesbox}

\end{document}
